\chapter{THỰC NGHIỆM, ĐÁNH GIÁ VÀ BÀN LUẬN KẾT QUẢ}

\section{Thiết lập thực nghiệm và hạ tầng phần cứng}

Các mô hình SLM của Phương pháp 1 (\texttt{Qwen3.5-2B} và \texttt{Qwen3.5-4B}) được huấn luyện trên Google Colab Pro với một GPU NVIDIA A100-SXM4-40GB (bộ nhớ khả dụng 39.49 GB), RAM 83.5 GB, Ubuntu 22.04 LTS, CUDA 12.4 và PyTorch 2.5. Nhóm sử dụng Unsloth [16] và QLoRA 4-bit NF4 [6]. Hai thành phần Bi-Encoder (\texttt{BGE-M3}) và Cross-Encoder (\texttt{xlm-roberta-base}) của Phương pháp 2 được huấn luyện độc lập trên GPU NVIDIA Tesla T4 16GB thuộc Kaggle. Do phần cứng và quy trình khác nhau, các số đo vận hành của hai phương pháp cần được diễn giải theo từng protocol.

Khi đánh giá, Method 1 chia dữ liệu thành hai phần và suy luận độc lập trên 2× NVIDIA Tesla T4 của Kaggle. Cách chia này rút ngắn thời gian chạy toàn tập, nhưng không giảm độ trễ của một truy vấn. Method 2 được đo trên 1× T4, batch size 1, có cache embedding công cụ và đồng bộ CUDA quanh từng truy vấn; phép đo Core dùng ba lượt warm-up cho mỗi tập.

Bảng 5.2 báo cáo độ trễ mean trên Core VI. Trong stress test, số đo của SLM là thời gian batch chia theo số mẫu, còn Method 2 báo cáo P50/P95 theo từng truy vấn. Các bước có ngẫu nhiên dùng seed 42; nghiên cứu chưa lặp nhiều seed để ước lượng độ lệch chuẩn hoặc khoảng tin cậy.

Đường dẫn tới mã nguồn, dữ liệu và các checkpoint được tập hợp tại Phụ lục A để đối chiếu từng cấu hình thực nghiệm.

\section{Hệ thống độ đo đánh giá (Evaluation Metrics)}

Tham khảo BFCL [4] và nghiên cứu của Ersoy và cộng sự [1], khóa luận sử dụng năm độ đo định lượng cho độ chính xác nội dung, khả năng từ chối, tính hợp lệ định dạng và độ trễ. Định nghĩa cụ thể dưới đây áp dụng cho protocol của khóa luận, không đồng nhất với toàn bộ thước đo ở các công trình nguồn.

Độ chính xác lựa chọn công cụ (Tool Selection Accuracy - $\text{Tool Acc}$, \%) phản ánh khả năng định tuyến đúng công cụ mục tiêu từ yêu cầu của người dùng. Chỉ số này được tính toán trên tập các câu truy vấn dương tính, đòi hỏi danh sách công cụ dự đoán phải trùng khớp chính xác với nhãn tham chiếu cả về định danh lẫn thứ tự xuất hiện đối với các trường hợp đa lệnh gọi (multi-call):
$$\text{Tool Acc} = \frac{N_{\text{tool-exact, positive}}}{N_{\text{positive}}}$$

Độ chính xác khớp toàn bộ đầu ra (Argument Population Accuracy - $\text{ArgA}$, \%) là tiêu chuẩn đánh giá khắt khe nhất trong bài toán gọi công cụ, đại diện cho tỷ lệ các trường hợp mô hình thực hiện thành công trọn vẹn tác vụ. Một dự đoán chỉ được công nhận là chính xác khi mô hình chọn đúng công cụ đồng thời trích xuất hoàn hảo toàn bộ các cặp khóa–giá trị tham số mà không có bất kỳ sai lệch nào sau bước chuẩn hóa. Chỉ số báo cáo mặc định trên toàn bộ tập kiểm thử ($\text{ArgA}_{\text{all}}$) bao gồm cả các trường hợp nhận diện chính xác câu truy vấn không gọi hàm:
$$\text{ArgA}_{\text{all}} = \frac{N_{\text{exact, all}}}{N_{\text{all}}}$$
Bên cạnh đó, khi cần phân tích chuyên biệt trên nhóm câu hỏi đòi hỏi kích hoạt công cụ, độ chính xác tham số trên tập dương tính ($\text{ArgA}_{\text{positive}}$) được xác định bởi:
$$\text{ArgA}_{\text{positive}} = \frac{N_{\text{exact, positive}}}{N_{\text{positive}}}$$
Trong toàn bộ khóa luận, ký hiệu $\text{ArgA}$ được quy ước đại diện cho $\text{ArgA}_{\text{all}}$ trừ khi có diễn giải riêng.

Độ thu hồi từ chối gọi hàm (Non-Function-Calling Recall - $\text{Non-FC Recall}$, \%) đo lường năng lực nhận diện các câu thoại giao tiếp thông thường không chứa ý định kích hoạt API. Đây là thước đo mang tính quyết định nhằm phát hiện và kiểm soát hiện tượng thiên kiến kích hoạt quá mức (over-triggering pathology) trong các hệ thống tác tử thực tế:
$$\text{Non-FC Recall} = \frac{N_{\text{true negative}}}{N_{\text{negative}}}$$

Tỷ lệ lỗi cú pháp (Syntax Error Rate, \%) xác định tỷ lệ đầu ra không thể được bộ phân tích cú pháp xử lý theo định dạng đánh giá. Với Method 2, đầu ra được kiến tạo từ schema nên protocol hiện tại không ghi nhận lỗi cú pháp; kết quả 0\% chỉ phản ánh tính hợp lệ về định dạng, không chứng tỏ chọn đúng công cụ hay điền đúng tham số và không thay thế chỉ số ArgA.

Cuối cùng, độ trễ suy luận (Wall-clock Latency, ms) lượng hóa thời gian từ khi tiếp nhận truy vấn đến khi tạo xong đầu ra. Khóa luận ghi rõ mean hoặc phân vị P50/P95, tập mẫu, phần cứng và phạm vi tính thời gian cho từng phép đo; những số liệu từ các protocol khác nhau chỉ dùng để mô tả, không tạo thành phép đo tốc độ đối chứng đồng nhất.

\section{Kết quả tổng thể và giải đáp 5 câu hỏi nghiên cứu}

Kết quả trên \texttt{CustomTools-VI} gồm 1,600 mẫu seen/unseen được trình bày ở Bảng 5.1a và 5.1b. Bảng 5.2 báo cáo kết quả trên 7,712 mẫu kiểm thử của mỗi ngôn ngữ trong \texttt{Canonical Core Benchmark}. Các bảng này là cơ sở phân tích năm câu hỏi nghiên cứu RQ1–RQ5.


\begin{table}[htbp]
    \centering
    \caption{a: Kết quả thực nghiệm trên tập kiểm thử CustomTools-VI Seen}
    \label{tab:bảng_5_1}
    \resizebox{\textwidth}{!}{%
    \begin{tabular}{llcccc}
        \toprule
        \textbf{Mô hình} & \textbf{Cấu hình huấn luyện} & \textbf{Tool Acc (\%)} & \textbf{ArgA (\%)} & \textbf{Non-FC Recall (\%)} & \textbf{Syntax Error (\%)} \\
        \midrule
        \textbf{Nhóm SLM cục bộ} &  &  &  &  &  \\
        \textbf{Qwen3.5-2B} & E0 (Zero-shot) & 39.25\% & 56.75\% & 97.75\% & 11.00\% \\
        \textbf{Qwen3.5-2B} & E1 (Đơn ngữ EN 60k) & 91.50\% & 63.12\% & 75.25\% & 4.50\% \\
        \textbf{Qwen3.5-2B} & E2 (Đơn ngữ VI 60k) & 91.00\% & 51.50\% & 67.50\% & 7.18\% \\
        \textbf{Qwen3.5-2B} & E3 (Song ngữ 60k) & 92.25\% & 19.38\% & 3.00\% & 4.31\% \\
        \textbf{Qwen3.5-2B} & E4 (Song ngữ + Miền VI) & 93.25\% & 87.00\% & \textbf{100.00\%} & 2.44\% \\
        \textbf{Qwen3.5-4B} & E3 (Song ngữ 60k) & 92.00\% & 62.00\% & 71.25\% & 12.00\% \\
        \textbf{Qwen3.5-4B} & \textbf{E4 (Song ngữ + Miền VI)} & \textbf{94.66\%} & \textbf{87.92\%} & \textbf{100.00\%} & 2.90\% \\
        \textbf{Nhóm Method 2 phân tách} &  &  &  &  &  \\
        \textbf{Method 2} & Bi-Encoder + Cross-Encoder & 92.25\% & 85.38\% & 92.75\% & \textbf{0.00\%} \\
        \textbf{Nhóm Frontier API đám mây} &  &  &  &  &  \\
        \textbf{GPT-5.6 Luna} & OpenAI API (không truyền tham số nhiệt độ) & 93.25\% & 78.25\% & 100.00\% & 0.00\% \\
        \textbf{Gemini 3.8 Flash} & Google API (không truyền tham số nhiệt độ) & \textbf{100.00\%} & \textbf{93.62\%} & 100.00\% & 0.06\% \\
        \bottomrule
    \end{tabular}
    }%
\end{table}

\textit{Ghi chú Bảng 5.1a: Các chỉ số được tính trên tập \texttt{test_seen} gồm 800 mẫu, với 400 mẫu dương tính và 400 mẫu âm tính. Tool Acc tính trên mẫu dương tính; ArgA tính trên toàn bộ mẫu; Non-FC Recall tính trên mẫu âm tính. Với Method 2, 0\% lỗi cú pháp là hệ quả của việc kiến tạo đầu ra theo schema, không đồng nghĩa với ArgA đạt 100\%. Mã benchmark API không truyền \texttt{temperature}; không thể quy kết kết quả cho cấu hình $T=0$.}


\begin{table}[htbp]
    \centering
    \caption{b: Kết quả thực nghiệm trên tập kiểm thử CustomTools-VI Unseen}
    \label{tab:bảng_5_1}
    \resizebox{\textwidth}{!}{%
    \begin{tabular}{llcccc}
        \toprule
        \textbf{Mô hình} & \textbf{Cấu hình huấn luyện} & \textbf{Tool Acc (\%)} & \textbf{ArgA (\%)} & \textbf{Non-FC Recall (\%)} & \textbf{Syntax Error (\%)} \\
        \midrule
        \textbf{Nhóm SLM cục bộ} &  &  &  &  &  \\
        \textbf{Qwen3.5-2B} & E0 (Zero-shot) & 40.75\% & 62.50\% & 97.25\% & 11.00\% \\
        \textbf{Qwen3.5-2B} & E1 (Đơn ngữ EN 60k) & 96.50\% & 70.50\% & 74.50\% & 4.50\% \\
        \textbf{Qwen3.5-2B} & E2 (Đơn ngữ VI 60k) & 96.25\% & 59.62\% & 67.75\% & 7.18\% \\
        \textbf{Qwen3.5-2B} & E3 (Song ngữ 60k) & 96.00\% & 27.88\% & 2.00\% & 4.31\% \\
        \textbf{Qwen3.5-2B} & E4 (Song ngữ + Miền VI) & 97.00\% & 86.38\% & \textbf{100.00\%} & 2.44\% \\
        \textbf{Qwen3.5-4B} & E3 (Song ngữ 60k) & 97.50\% & 69.75\% & 74.50\% & 12.00\% \\
        \textbf{Qwen3.5-4B} & \textbf{E4 (Song ngữ + Miền VI)} & 96.75\% & \textbf{86.75\%} & \textbf{100.00\%} & 2.90\% \\
        \textbf{Nhóm Method 2 phân tách} &  &  &  &  &  \\
        \textbf{Method 2} & Bi-Encoder + Cross-Encoder & 79.50\% & 60.25\% & 93.75\% & \textbf{0.00\%} \\
        \textbf{Nhóm Frontier API đám mây} &  &  &  &  &  \\
        \textbf{GPT-5.6 Luna} & OpenAI API (không truyền tham số nhiệt độ) & 97.25\% & 79.50\% & 100.00\% & 0.00\% \\
        \textbf{Gemini 3.8 Flash} & Google API (không truyền tham số nhiệt độ) & \textbf{100.00\%} & \textbf{92.12\%} & 99.75\% & 0.06\% \\
        \bottomrule
    \end{tabular}
    }%
\end{table}

\textit{Ghi chú Bảng 5.1b: Các chỉ số được tính trên tập \texttt{test_unseen} gồm 800 mẫu, với 400 mẫu dương tính và 400 mẫu âm tính. Tool Acc tính trên mẫu dương tính; ArgA tính trên toàn bộ mẫu; Non-FC Recall tính trên mẫu âm tính. Với Method 2, 0\% lỗi cú pháp là hệ quả của việc kiến tạo đầu ra theo schema, không đồng nghĩa với ArgA đạt 100\%. Mã benchmark API không truyền \texttt{temperature}; không thể quy kết kết quả cho cấu hình $T=0$.}

\begin{figure}[htbp]
    \centering
    \includegraphics[width=0.88\textwidth]{figures/fig2_performance_comparison}
    \caption{So sánh hiệu năng các phương pháp trên benchmark CustomTools-VI}
    \label{fig:hình_5_1}
\end{figure}

\textit{Hình 5.1: So sánh hiệu năng các phương pháp trên benchmark CustomTools-VI (Seen vs. Unseen).}



\begin{table}[htbp]
    \centering
    \caption{Kết quả thực nghiệm trên Canonical Core Benchmark (7,712 mẫu)}
    \label{tab:bảng_5_2}
    \resizebox{\textwidth}{!}{%
    \begin{tabular}{llccccccc}
        \toprule
        \textbf{Mô hình} & \textbf{Cấu hình huấn luyện} & \textbf{VI Test: Tool Acc (\%)} & \textbf{VI Test: ArgA (\%)} & \textbf{VI Test: Non-FC (\%)} & \textbf{EN Test: Tool Acc (\%)} & \textbf{EN Test: ArgA (\%)} & \textbf{EN Test: Non-FC (\%)} & \textbf{Độ trễ VI mean (ms)} \\
        \midrule
        \textbf{Qwen3.5-2B} & E0 (Zero-Shot) & 56.34\% & 40.48\% & 96.33\% & 85.18\% & 60.63\% & 94.50\% & 707 ms \\
        \textbf{Qwen3.5-2B} & E1 (Đơn ngữ EN 60k) & 93.67\% & 65.57\% & 94.30\% & 94.07\% & \textbf{73.66\%} & 94.30\% & 878 ms \\
        \textbf{Qwen3.5-2B} & E2 (Đơn ngữ VI 60k) & 90.25\% & 64.90\% & 94.30\% & 93.45\% & 71.36\% & 93.08\% & 872 ms \\
        \textbf{Qwen3.5-2B} & E3 (Song ngữ 60k) & 94.00\% & 69.76\% & 94.30\% & 94.27\% & 73.22\% & 94.30\% & 864 ms \\
        \textbf{Qwen3.5-2B} & E4 (Song ngữ + Miền VI) & 93.92\% & 69.75\% & 94.30\% & 94.17\% & 73.15\% & 94.30\% & 970 ms \\
        \textbf{Qwen3.5-4B} & E3 (Song ngữ 60k) & \textbf{98.73\%} & \textbf{72.86\%} & 94.30\% & \textbf{96.63\%} & \textbf{74.71\%} & 94.30\% & 2,432 ms \\
        \textbf{Qwen3.5-4B} & E4 (Song ngữ + Miền VI) & 86.22\% & 64.94\% & 94.30\% & 85.03\% & 66.66\% & 94.30\% & 2,485 ms \\
        \textbf{Method 2} & Shared E4 (Bi+Cross) & 59.20\% & 30.26\% & 94.09\% & 62.60\% & 35.52\% & 94.30\% & \textbf{61.50 ms} \\
        \bottomrule
    \end{tabular}
    }%
\end{table}

\textit{Ghi chú: Cột độ trễ báo cáo mean trên Core VI. SLM chia dữ liệu trên 2× T4 và suy luận độc lập từng phần, không tăng tốc một truy vấn; Method 2 đo trên 1× T4, batch 1, có cache embedding công cụ và đồng bộ CUDA quanh từng truy vấn (P50 = 59.55 ms). Method 2 được đánh giá trên đủ 7,712 mẫu của mỗi ngôn ngữ. Do phạm vi xử lý và protocol đo khác nhau, không tính tỷ lệ tăng tốc trực tiếp từ cột này.}


\section{1. RQ1: Năng lực chuyển giao tri thức ngôn ngữ chéo (Cross-Lingual Transfer)}

Trên Core VI (Bảng 5.2), E1 chỉ huấn luyện bằng tiếng Anh đạt \textbf{65.57\%} ArgA, cao hơn E2 chỉ huấn luyện bằng tiếng Việt (\textbf{64.90\%}). Trên CustomTools-VI Unseen, hai cấu hình lần lượt đạt \textbf{70.50\%} và \textbf{59.62\%}. Những kết quả này cho thấy E1 vẫn xử lý được truy vấn tiếng Việt trong hai tập đánh giá.

Hiện tượng thực nghiệm này phù hợp với giả thuyết rằng Qwen3.5 chuyển giao được một phần năng lực phân tích schema từ tiếng Anh sang tiếng Việt. Tuy nhiên, thiết kế hiện tại chưa tách riêng tác động của tiền huấn luyện đa ngôn ngữ và dữ liệu tinh chỉnh. Ở chiều đánh giá ngược trên Core EN, E2 đạt 71.36\% so với 73.66\% của E1; E3 đạt ArgA 69.76\% trên Core VI ở mô hình 2B và 72.86\% ở mô hình 4B. Các kết quả cho thấy sự khác biệt giữa cấu hình dữ liệu trong những lần chạy này, chưa chứng minh E3 giải quyết hoàn toàn mất cân bằng ngôn ngữ.

\section{2. RQ2: Hiện tượng kích hoạt công cụ quá mức và vai trò của dữ liệu âm tính}

Trên CustomTools-VI, E3 được huấn luyện bằng 60,000 mẫu song ngữ nhưng không có dữ liệu âm tính đặc thù miền. Non-FC Recall chỉ đạt \textbf{2.00\%–3.00\%}, nghĩa là mô hình gọi công cụ nhầm ở hơn 97\% truy vấn âm tính trong hai tập kiểm thử. Đây là biểu hiện của thiên kiến kích hoạt công cụ quá mức (over-triggering) trong điều kiện đánh giá này.

Trên CustomTools-VI, ArgA-all của SLM 2B E3 đạt \textbf{19.38\%} ở seen và \textbf{27.88\%} ở unseen. Sau khi bổ sung 5,600 mẫu bản địa ở E4, trong đó có mẫu âm tính, Non-FC Recall đạt \textbf{100.00\%} trên cả hai tập kiểm thử; ArgA-all tăng lên \textbf{87.00\%} và \textbf{86.38\%}. Chênh lệch phù hợp với giả thuyết dữ liệu bản địa hỗ trợ nhận diện truy vấn không gọi hàm, nhưng E3 và E4 khác nhau ở nhiều khía cạnh dữ liệu nên chưa xác định riêng tác động của mẫu âm tính.

\section{3. RQ3: Khả năng tổng quát hóa Zero-Shot trên công cụ chưa từng gặp}

Trên \texttt{test_unseen}, SLM E4 và Method 2 có mức suy giảm ArgA khác nhau so với \texttt{test_seen}.

SLM E4 đạt ArgA \textbf{86.38\%} ở bản 2B và \textbf{86.75\%} ở bản 4B trên tập unseen, chênh lần lượt $-0.62$ và $-1.17$ điểm phần trăm so với tập seen. Kết quả phù hợp với khả năng sử dụng mô tả công cụ và schema trong prompt để xử lý công cụ mới; phép đánh giá hiện tại chưa tách riêng tác động của từng thành phần này.

Method 2 đạt ArgA \textbf{85.38\%} trên \texttt{test_seen} và \textbf{60.25\%} trên \texttt{test_unseen}; ArgA của riêng truy vấn dương tính unseen là 26.75\%. Tool Acc giảm từ 92.25\% xuống 79.50\%. Các số liệu cho thấy cả việc chọn công cụ lẫn hoàn thiện lời gọi đều khó hơn trên tập unseen. Đối chứng Oracle ở Chương 6 giúp phân tích thêm vai trò của từng giai đoạn, nhưng chưa xác định nguyên nhân duy nhất của mức suy giảm.

\section{4. RQ4: Đánh đổi Pareto giữa độ chính xác, độ trễ và chi phí tính toán}

Nhằm cung cấp cơ sở khoa học phục vụ việc lựa chọn kiến trúc trong các kịch bản triển khai thực tế, Bảng 5.3 tổng hợp đường biên đánh đổi đa chiều giữa các phương pháp đại diện trên toàn bộ các trục tiêu chí kỹ thuật cốt lõi.


\begin{table}[htbp]
    \centering
    \caption{Đối chiếu chỉ số và điều kiện vận hành của bốn phương pháp}
    \label{tab:bảng_5_3}
    \resizebox{\textwidth}{!}{%
    \begin{tabular}{lccccl}
        \toprule
        \textbf{Trục đánh giá kỹ thuật} & \textbf{Method 1: SLM Qwen3.5-2B (E4)} & \textbf{Method 1: SLM Qwen3.5-4B (E4)} & \textbf{Method 2: Bi+Cross (BGE-M3 + XLM-R)} & \textbf{Frontier Closed: GPT-5.6 Luna} & \textbf{Khuyến nghị lựa chọn kỹ thuật} \\
        \midrule
        \textbf{Kiến trúc mô hình} & Tự hồi quy (Decoder-only) & Tự hồi quy (Decoder-only) & Phân tách phân biệt (Encoders) & Mô hình đám mây (Cloud API) & Tùy biến theo năng lực hạ tầng \\
        \textbf{Quy mô tham số} & 2.22 tỷ (QLoRA 10.9M) & 4.56 tỷ (QLoRA 21.2M) & 567M (Bi) + 278M (Cross) & Không công bố trong log đánh giá & Phụ thuộc kiến trúc và triển khai \\
        \textbf{Custom Seen ArgA} & 87.00\% & \textbf{87.92\%} & 85.38\% & 78.25\% & SLM E4 cao hơn Luna trên tập này \\
        \textbf{Custom Unseen ArgA} & 86.38\% & \textbf{86.75\%} & 60.25\% & 79.50\% & SLM E4 cao hơn trên tập này \\
        \textbf{Lỗi cú pháp CustomTools-VI} & 2.44\% & 2.90\% & \textbf{0.00\% (đầu ra kiến tạo từ schema)} & 0.00\% & Chỉ phản ánh định dạng; xem ArgA để đánh giá nội dung \\
        \textbf{Độ trễ Core VI mean (ms)\}* & 970 ms & 2,485 ms & \textbf{61.50 ms} & Không đo trên Core VI & Chỉ đối chiếu mô tả do protocol khác nhau \\
        \textbf{Mức tiêu thụ VRAM (đo đạc)\}* & ~4.8–5.5 GB VRAM & ~9.5 GB VRAM & \textbf{3,283 MiB allocated; 3,772 MiB reserved} & Không tiêu tốn VRAM cục bộ & Số liệu phụ thuộc protocol và công cụ đo \\
        \textbf{Khả năng mở rộng danh mục} & Giới hạn bởi cửa sổ ngữ cảnh & Giới hạn bởi cửa sổ ngữ cảnh & Đã kiểm thử với $N \le 1000$ & Chưa kiểm thử stress test & Cần đánh giá theo quy mô danh mục thực tế \\
        \textbf{Phạm vi quan sát} & Core VI và CustomTools-VI & Core VI và CustomTools-VI & Core VI, CustomTools-VI và stress test & CustomTools-VI & Chọn theo miền dữ liệu và điều kiện triển khai \\
        \bottomrule
    \end{tabular}
    }%
\end{table}

\textit{Ghi chú: SLM dùng hai T4 để chia dữ liệu và suy luận độc lập; một truy vấn chạy trên một GPU. Method 2 dùng một T4, batch 1 và cache embedding. Các số trong hàng độ trễ là mean Core VI của cấu hình E4, không phải P50; P50 từng truy vấn của Method 2 trong stress test nằm ở Bảng 5.4b. Hàng VRAM kết hợp số đo từ các run và cách ghi nhận khác nhau; riêng Method 2 là bộ nhớ PyTorch allocated/reserved tại $N=1000$. Protocol đo và phạm vi xử lý khác nhau nên các hàng này không chứng minh tỷ lệ tăng tốc hay mức tiết kiệm VRAM trực tiếp. API chỉ được kiểm thử trên CustomTools-VI.}

Các kết quả gợi ý hai lựa chọn kiến trúc theo yêu cầu triển khai. Trên CustomTools-VI unseen, SLM E4 đạt ArgA cao hơn Method 2, phù hợp với bài toán thường xuyên tiếp nhận schema mới. Trong protocol stress test với danh mục xác định trước, Method 2 ghi nhận P50 dưới 100 ms ở $N \le 500$, tiếp tục xử lý được $N=1000$ và không có lỗi cú pháp theo định nghĩa đầu ra kiến tạo từ schema. Khả năng áp dụng cho thiết bị biên và mục tiêu độ trễ sản xuất vẫn cần đo trên phần cứng, tải và quy trình phục vụ tương ứng.

\section{5. RQ5: Kiểm thử áp lực quy mô lớn và giới hạn vật lý của mô hình tạo sinh}

Để kiểm chứng tính bền bỉ của hai trường phái kiến trúc khi quy mô danh mục công cụ mở rộng từ dải nhỏ ($N = 3$) đến quy mô thực tế của các nền tảng doanh nghiệp ($N = 1000$ công cụ), nghiên cứu tiến hành bài kiểm thử áp lực (\textbf{Stress Test}) đối đầu trực diện giữa đại diện tiêu biểu của Method 1 (\texttt{Qwen3.5-2B} E4) và Method 2 trên cùng 200 truy vấn Custom Seen (1,200 mẫu thử nghiệm lồng nhau). Độ chính xác được báo cáo ở Bảng 5.4a, độ trễ và bộ nhớ ở Bảng 5.4b; Hình 5.2 trực quan hóa xu hướng theo $N$.


\begin{table}[htbp]
    \centering
    \caption{a: Độ chính xác trong Stress Test theo quy mô danh mục công cụ ($N = 3 \to 1000$)}
    \label{tab:bảng_5_4}
    \begin{tabular}{crrrr}
        \toprule
        \textbf{Số công cụ ($N$)} & \textbf{SLM Tool Acc (\%)} & \textbf{M2 Tool Acc (\%)} & \textbf{SLM ArgA (\%)} & \textbf{M2 ArgA (\%)} \\
        \midrule
        \textbf{3} & \textbf{94.0\%} & 89.0\% & 85.5\% & \textbf{87.5\%} \\
        \textbf{10} & \textbf{93.0\%} & 89.0\% & 84.5\% & \textbf{87.5\%} \\
        \textbf{50} & \textbf{93.0\%} & 89.0\% & 85.5\% & \textbf{87.5\%} \\
        \textbf{100} & \textbf{92.0\%} & 89.0\% & 83.0\% & \textbf{87.0\%} \\
        \textbf{500} & \textit{OOM} & \textbf{87.0\%} & \textit{OOM} & \textbf{85.0\%} \\
        \textbf{1000} & \textit{OOM} & \textbf{87.0\%} & \textit{OOM} & \textbf{84.0\%} \\
        \bottomrule
    \end{tabular}
\end{table}


\begin{table}[htbp]
    \centering
    \caption{b: Độ trễ và bộ nhớ trong Stress Test theo quy mô danh mục công cụ}
    \label{tab:bảng_5_4}
    \begin{tabular}{crrr}
        \toprule
        \textbf{Số công cụ ($N$)} & \textbf{Độ trễ SLM (ms)*} & \textbf{M2 P50 (ms)} & \textbf{M2 P95 (ms)} \\
        \midrule
        \textbf{3} & 1,258.94 & \textbf{55.13} & 84.92 \\
        \textbf{10} & 1,604.68 & \textbf{55.05} & 84.74 \\
        \textbf{50} & 4,703.39 & \textbf{56.07} & 86.46 \\
        \textbf{100} & 9,318.21 & \textbf{61.01} & 85.09 \\
        \textbf{500} & \textit{OOM} & \textbf{92.27} & 128.82 \\
        \textbf{1000} & \textit{OOM} & \textbf{107.68} & 148.21 \\
        \bottomrule
    \end{tabular}
\end{table}

\textit{Ghi chú Bảng 5.4b: Độ trễ SLM là thời gian sinh của batch chia cho số mẫu trong batch, với batch size thay đổi theo $N$; đây không phải P50 wall-clock từng truy vấn. Method 2 đo từng truy vấn ở batch 1, có đồng bộ CUDA và cache embedding. Hai phép đo khác giao thức, nên không quy đổi thành hệ số tăng tốc. PyTorch allocated VRAM của Method 2 tại $N=3, 10, 50, 100, 500, 1000$ lần lượt là 3,277; 3,282; 3,278; 3,283; 3,283; 3,283 MiB.}

\begin{figure}[htbp]
    \centering
    \includegraphics[width=0.88\textwidth]{figures/fig3_stress_test_curves}
    \caption{Kết quả Stress Test đối đầu trực diện khi mở rộng quy mô công cụ ($N = 3 \to 1000$)}
    \label{fig:hình_5_2}
\end{figure}

\textit{Hình 5.2: Kết quả Stress Test đối đầu trực diện ($N = 3 \to 1000$): (a) Độ chính xác và ranh giới CUDA OOM của SLM tại $N \ge 500$; (b) Độ trễ suy luận trên thang log (P50 đơn truy vấn ở Method 2 và thời gian batch chia đều ở SLM).}

Thực nghiệm kiểm thử áp lực đã làm sáng tỏ những khác biệt mang tính bản chất trong cơ chế vận hành của hai mô thức thiết kế:

Trong stress test giữ nguyên tập truy vấn mỏ neo, ArgA của Method 2 giảm từ 87.50\% ở $N=3$ xuống 84.00\% ở $N=1000$, tức 3.50 điểm phần trăm. Ở $N \le 100$, ArgA của Method 2 cao hơn SLM từ 2.0 đến 4.0 điểm phần trăm; tại $N=1000$, Non-FC Recall của Method 2 là 95.00\%. Quan sát cho thấy cơ chế lọc ngưỡng ($\tau=0.35$, $\delta=0.21$) vẫn hoạt động trong protocol này, nhưng chưa tách riêng đóng góp của nó khỏi các thành phần khác.

Trong protocol stress test, P50 từng truy vấn của Method 2 nằm trong khoảng 55.05–61.01 ms ở $N \le 100$ và đạt 107.68 ms tại $N=1000$. Việc phải đối chiếu với nhiều vector công cụ hơn có thể góp phần làm tăng thời gian xử lý. Số đo SLM từ thời gian batch chia theo mẫu tăng từ 1,258.94 ms ($N=3$) lên 9,318.21 ms ($N=100$), phù hợp với việc prompt chứa thêm schema; hai chuỗi số liệu không cùng định nghĩa nên không dùng để tính hệ số tăng tốc.

Trong cấu hình stress test trên GPU 16GB, SLM gặp CUDA Out of Memory ở $N \ge 500$. Quan sát này phù hợp với giả thuyết rằng prompt dài làm tăng nhu cầu bộ nhớ ở các lớp attention toàn phần, song chưa đủ để tách riêng tác động của attention, KV cache, precision, batch và cách hiện thực kernel. Cấu hình Qwen3.5 được dùng còn xen kẽ các lớp linear attention và full attention [17], nên không thể dùng riêng độ phức tạp $\mathcal{O}(L^2)$ của full attention để giải thích toàn bộ hiện tượng. Method 2 tiếp tục chạy ở $N=1000$ và ghi nhận PyTorch allocated VRAM khoảng 3.21 GiB (3,283 MiB) trong protocol này.

\section{Nghiên cứu mở rộng quy mô mô hình: Qwen3.5-2B vs. Qwen3.5-4B}

Để làm sáng tỏ ảnh hưởng của dung lượng tham số đối với năng lực tiếp nhận schema phức tạp và tính ổn định cú pháp, quá trình huấn luyện được mở rộng lên mô hình \texttt{Qwen3.5-4B} (4.56 tỷ tham số). Kết quả đối chuẩn chi tiết được trình bày tại Bảng 5.5, Bảng 5.6 và minh họa tại Hình 5.3.


\begin{table}[htbp]
    \centering
    \caption{Nghiên cứu mở rộng quy mô tham số mô hình (Qwen3.5-2B vs. 4B)}
    \label{tab:bảng_5_5}
    \resizebox{\textwidth}{!}{%
    \begin{tabular}{llcccccc}
        \toprule
        \textbf{Cấu hình} & \textbf{Backbone} & \textbf{Core VI ArgA (\%)} & \textbf{Custom Seen ArgA (\%)} & \textbf{Custom Unseen ArgA (\%)} & \textbf{ArgA Gap} & \textbf{Cú pháp lỗi Core VI (\%)} & \textbf{Độ trễ Core VI mean (ms)} \\
        \midrule
        \textbf{E3 (Song ngữ)} & Qwen3.5-2B & 69.76\% & 19.38\% & 27.88\% & +8.50\% & 4.31\% & 864 ms \\
        \textbf{E3 (Song ngữ)} & \textbf{Qwen3.5-4B} & \textbf{72.86\%} & \textbf{62.00\%} & \textbf{69.75\%} & +7.75\% & \textbf{0.32\%} & 2,432 ms \\
        \textbf{E4 (Đặc thù miền)} & Qwen3.5-2B & 69.75\% & 87.00\% & 86.38\% & -0.62\% & 4.67\% & 970 ms \\
        \textbf{E4 (Đặc thù miền)} & \textbf{Qwen3.5-4B} & 64.94\% & \textbf{87.92\%} & \textbf{86.75\%} & \textbf{-1.17\%} & 9.75\% & 2,485 ms \\
        \bottomrule
    \end{tabular}
    }%
\end{table}


\begin{table}[htbp]
    \centering
    \caption{Hiệu năng đối đầu của mô hình Qwen3.5-4B (E3 vs. E4) trên Core Benchmark}
    \label{tab:bảng_5_6}
    \resizebox{\textwidth}{!}{%
    \begin{tabular}{llcccccc}
        \toprule
        \textbf{Cấu hình} & \textbf{Tập kiểm thử} & \textbf{Số mẫu (Samples)} & \textbf{Tool Selection Acc (\%)} & \textbf{ArgA / Exact Match (\%)} & \textbf{Non-FC Recall (\%)} & \textbf{Tỷ lệ lỗi cú pháp (\%)} & \textbf{Độ trễ trung bình (ms)} \\
        \midrule
        \textbf{E3 (Song ngữ)} & Core VI Test & 7,712 (7,221 pos / 491 neg) & \textbf{98.73\%} & \textbf{72.86\%} & 94.30\% & \textbf{0.32\%} & 2,431.83 ms \\
        \textbf{E3 (Song ngữ)} & Core EN Test & 7,712 (7,221 pos / 491 neg) & \textbf{96.63\%} & \textbf{74.71\%} & 94.30\% & \textbf{1.97\%} & 3,289.59 ms \\
        \textbf{E4 (Đặc thù miền)} & Core VI Test & 7,712 (7,221 pos / 491 neg) & 86.22\% & 64.94\% & 94.30\% & 9.75\% & 2,485.12 ms \\
        \textbf{E4 (Đặc thù miền)} & Core EN Test & 7,712 (7,221 pos / 491 neg) & 85.03\% & 66.66\% & 94.30\% & 11.36\% & 3,310.25 ms \\
        \bottomrule
    \end{tabular}
    }%
\end{table}

\begin{figure}[htbp]
    \centering
    \includegraphics[width=0.88\textwidth]{figures/fig4_model_scaling}
    \caption{Nghiên cứu mở rộng quy mô tham số mô hình SLM (2B vs. 4B)}
    \label{fig:hình_5_3}
\end{figure}

\textit{Hình 5.3: So sánh mô hình 2B và 4B: (a) Độ chính xác chọn công cụ và trích xuất tham số trên Core Benchmark; (b) Tỷ lệ lỗi cú pháp JSON.}

Mục này so sánh kết quả của hai checkpoint Qwen3.5 có quy mô 2.22 tỷ và 4.56 tỷ tham số trong các cấu hình đã thử nghiệm.

Ở cấu hình song ngữ E3 trên Core VI, tỷ lệ vi phạm cú pháp giảm từ 4.31\% ở mô hình 2B xuống \textbf{0.32\%} ở mô hình 4B, tương đương mức giảm hơn 13 lần trong phép đo này. Chênh lệch phù hợp với giả thuyết mô hình lớn hơn tuân thủ định dạng tốt hơn, nhưng chưa chứng minh riêng quy mô tham số là nguyên nhân vì hai checkpoint có thể khác nhau ở nhiều yếu tố.

Trên Core VI, Qwen3.5-4B E3 đạt Tool Acc \textbf{98.73\%} và ArgA \textbf{72.86\%}, cao hơn bản 2B E3 lần lượt 4.73 và 3.10 điểm phần trăm. Phép so sánh này ghi nhận lợi thế của checkpoint 4B trong cấu hình E3, nhưng chưa tách riêng ảnh hưởng của quy mô tham số khỏi các khác biệt huấn luyện khác.

Trên CustomTools-VI, ArgA của 2B E3 là 19.38\% ở seen và 27.88\% ở unseen; 4B E3 lần lượt đạt 62.00\% và 69.75\%. Khoảng cách này cho thấy hai checkpoint phản ứng khác nhau với dữ liệu miền mới, song chưa đủ để quy toàn bộ khác biệt cho kích thước mô hình.

Sau khi bổ sung 5,600 mẫu nghiệp vụ Việt Nam ở E4, Qwen3.5-4B đạt ArgA \textbf{87.92\%} trên seen và \textbf{86.75\%} trên unseen của CustomTools-VI. Trên Core VI, ArgA của E4 còn 64.94\% và tỷ lệ lỗi cú pháp là 9.75\%. Vì vậy, lựa chọn giữa E3 và E4 cần dựa trên miền dữ liệu mục tiêu; kết quả hiện có không bảo đảm checkpoint nào tốt hơn trong mọi điều kiện triển khai.

\section{So sánh đối đầu với các Frontier API thương mại đóng (GPT-5.6 Luna, Gemini 3.8 Flash)}

Nghiên cứu đánh giá các mô hình đề xuất cùng hai mô hình API là GPT-5.6 Luna và Gemini 3.8 Flash trên 1,600 mẫu CustomTools-VI. Benchmark dùng cùng câu lệnh hệ thống, truy vấn và danh mục JSON Schema cho hai API qua Kaggle Benchmark SDK. Mã đánh giá không truyền tham số \texttt{temperature}, vì vậy không xác nhận được nhiệt độ thực tế hoặc loại trừ biến thiên ngẫu nhiên của dịch vụ [18], [19].

Trên CustomTools-VI (Bảng 5.1a và Bảng 5.1b), Qwen3.5-4B E4 đạt ArgA \textbf{87.92\%} trên seen và \textbf{86.75\%} trên unseen, cao hơn GPT-5.6 Luna (78.25\% và 79.50\%) lần lượt \textbf{9.67} và \textbf{7.25} điểm phần trăm. Qwen3.5-2B E4 cũng cao hơn lần lượt 8.75 và 6.88 điểm phần trăm. Kết quả gợi ý lợi ích của tinh chỉnh theo miền trên bộ dữ liệu này; một benchmark và một lần chạy không đủ để quy toàn bộ chênh lệch cho fine-tuning hoặc khái quát sang các miền khác.

Gemini 3.8 Flash đạt ArgA khoảng 92\%–93\% trên hai tập kiểm thử CustomTools-VI trong lần chạy này. Ngoài độ chính xác, việc lựa chọn giữa mô hình cục bộ và dịch vụ API cần xét thêm quản trị dữ liệu, chi phí và năng lực vận hành:

Triển khai cục bộ có thể giảm phụ thuộc vào việc truyền dữ liệu nghiệp vụ đến nhà cung cấp bên ngoài và tạo điều kiện kiểm soát luồng dữ liệu. Tuy nhiên, việc tuân thủ pháp luật bảo vệ dữ liệu cá nhân và an ninh mạng vẫn phụ thuộc vào mục đích xử lý, phân quyền, lưu log và vận hành hạ tầng; mô hình cục bộ không loại bỏ nguy cơ rò rỉ. Tại thời điểm khóa luận, các văn bản cần xem xét bao gồm Luật Bảo vệ dữ liệu cá nhân số 91/2025/QH15 và Luật An ninh mạng số 116/2025/QH15 [20], [21].

Chi phí triển khai cục bộ gồm đầu tư phần cứng và chi phí vận hành, bảo trì, điện năng; chi phí API thường phụ thuộc vào mức sử dụng. Khóa luận chưa đo tổng chi phí sở hữu theo lưu lượng nên không kết luận phương án nào rẻ hơn trong dài hạn.

Triển khai cục bộ cũng cho phép chủ động quản lý đường ống suy luận. Mức độ sẵn sàng thực tế vẫn phụ thuộc vào năng lực vận hành, phần cứng và cơ chế dự phòng của từng tổ chức.

\section{1. Thông tin tái lập phép đánh giá Frontier API}

Hai run đầy đủ trên CustomTools-VI bắt đầu ngày 16/09/2026 (UTC), sử dụng task Kaggle Benchmark \texttt{tool-calling-vi-full} phiên bản 1 và agent \texttt{kaggle-benchmarks} 0.6.1. Model slug trong log lần lượt là \texttt{openai/gpt-5.6-luna} và \texttt{google/gemini-3.8-flash}; model ID tương ứng được nhà cung cấp công bố là \texttt{gpt-5.6-luna} [18] và \texttt{gemini-3.8-flash} [19]. Mã task gọi giao diện \texttt{llm.client.chat.completions.create}, truyền \texttt{model=llm.model}, \texttt{tool_choice="auto"} và danh sách công cụ lấy từ \texttt{tools} của từng mẫu. Riêng GPT-5.6 Luna được truyền \texttt{reasoning_effort="none"}. Tám worker xử lý các mẫu song song, nên độ trễ API nếu được báo cáo là thời gian wall-clock từng lời gọi trong điều kiện có đồng thời nhiều request, bao gồm cả mạng và dịch vụ.

Prompt hệ thống dùng cho cả hai API là: “Bạn là trợ lý AI hữu ích. Nếu người dùng yêu cầu hành động cần công cụ, hãy gọi đúng các hàm và truyền tham số chính xác. Nếu yêu cầu không liên quan đến công cụ nào được cung cấp, TUYỆT ĐỐI KHÔNG gọi công cụ nào.” Prompt người dùng lấy từ \texttt{query}; định nghĩa function gồm \texttt{name}, \texttt{description} và \texttt{parameters} được chuyển từ trường \texttt{tools}.

Mã task không truyền \texttt{temperature}, giới hạn output token hoặc tham số retry tường minh. Log có model slug và thời điểm chạy, nhưng không có snapshot dịch vụ, endpoint thực tế, region hoặc chính sách retry của SDK. Vì vậy, các điều kiện này có thể khác khi chạy lại.

